% !TeX root = Chapter_03_Slides.tex
% Chapter 3 lecture slides — Risk Identification: Categorizing and Discovering Enterprise Risks
% Instructor-resource series.
\documentclass[11pt, aspectratio=169]{beamer}
\usepackage{tikz}
\makeatletter
\def\input@path{{./}{../slides/}}
\makeatother
\usepackage{erm_beamer_theme}
\usepackage{amsmath,amssymb}

\title[ERM Ch.\,3]{Risk Identification: Categorizing and Discovering Enterprise Risks}
\subtitle{Enterprise Risk Management, Second Edition --- Chapter 3}
\author{Martin F.\ Grace}
\institute{University of Iowa\\[2pt] Tippie College of Business\\[2pt] Vaughan Institute of Risk Management}
\date{June 2026}

\begin{document}

\begin{frame}[plain]
\titlepage
\end{frame}

\begin{frame}{Learning Objectives}
\begin{itemize}
  \item Explain why \alert{risk identification} is the foundational step in ERM.
  \item Distinguish among hazard, financial, operational, and strategic risks.
  \item Explain why financial \emph{loss} is not the same as financial \emph{risk}, and classify business interruption correctly.
  \item Describe how risks interact across quadrant boundaries and why siloed identification fails.
  \item Apply internal and external identification methodologies to a real organization.
  \item Read and critically evaluate the \alert{Risk Factors} section of an SEC Form~10-K.
\end{itemize}
\end{frame}

\section{Opening Case: Netflix and the Red Envelope Problem}

\begin{frame}{The Red Envelope Business}
\begin{itemize}
  \item Early 2000s: DVD-by-mail \emph{was} the business --- millions of subscribers, growing revenue, a sophisticated logistics operation.
  \item Blockbuster's brick-and-mortar stores suddenly looked slow and expensive. Wall Street liked the story.
  \item Inside the company, executives and engineers were tracking something with nothing to do with postal routes: \alert{broadband penetration} and rapidly improving video compression.
  \item If movies could be delivered instantly over the internet, what happens to a business built on mailing plastic discs?
\end{itemize}
\end{frame}

\begin{frame}{Naming the Risk --- in Writing}
\begin{itemize}
  \item Netflix named the obsolescence of its own business model as a concrete, enterprise-level \alert{strategic risk} --- in board discussions and in the Risk Factors of its Form~10-K.
  \item Disclosed uncertainties, stated \emph{while the DVD business still looked strong}:
  \begin{itemize}
    \item Broadband might spread slower than expected
    \item Studios might restrict or overprice digital rights
    \item Competitors with ``greater financial, marketing, and technical resources'' could win streaming
    \item The physical-to-digital shift could alienate customers
  \end{itemize}
  \item Acting on it meant spending before the payoff was certain: streaming engineers, content licenses, ``Watch Instantly'' --- all funded by red-envelope cash flow.
\end{itemize}
\end{frame}

\begin{frame}{The Lesson: Identification, Not Just Innovation}
\begin{itemize}
  \item The transition was messy --- the 2011 brand-split debacle tanked the stock. But the underlying risk assessment never changed.
  \item A decade later: red envelopes gone, Blockbuster gone, Netflix a global streaming platform.
  \item The most important risks don't always arrive as sudden crises. Some arrive as \alert{slow-moving shifts} that turn yesterday's strengths into tomorrow's vulnerabilities.
\end{itemize}
\medskip
\begin{block}{The identification problem}
Netflix saw that its successful business model contained built-in obsolescence risk --- and wrote it down before the old model had obviously failed. Why is that so hard? What would it take for \emph{your} organization to disclose that its core business is the risk?
\end{block}
\end{frame}

\section{The Four Quadrants of Risk}

\begin{frame}{Source, Not Consequence}
\begin{itemize}
  \item The taxonomy: \alert{hazard}, \alert{financial}, \alert{operational}, \alert{strategic}. Each quadrant has a distinct source, characteristic tools, and a different relationship to strategy.
  \item Almost every risk ultimately produces financial consequences --- so consequence \emph{cannot} be the basis for classification.
\end{itemize}
\medskip
\begin{block}{The classification question}
Not ``does this cost money?'' but \alert{``what causes this uncertainty?''} A warehouse fire is hazard risk (physical peril); a euro drop is financial risk (market prices). Getting the source right determines which tools apply and who is responsible.
\end{block}
\end{frame}

\begin{frame}{Hazard Risks: Pure Downside}
\begin{itemize}
  \item \alert{Pure risks}: loss or no loss --- no upside. The historical territory of corporate risk management.
  \item \textbf{Property and natural disaster}: fire, flood, earthquake, windstorm; climate change is shifting frequency and severity faster than historical actuarial data can track.
  \item \textbf{Liability}: product, general, professional, employment practices, D\&O, and --- fastest growing --- \alert{cyber liability} (nine or ten figures for large breaches).
  \item Four management strategies, combined by cost and risk character:
  \begin{itemize}
    \item \textbf{Control} (sprinklers, training) --- \textbf{Insurance} (pooling makes these risks transferable) --- \textbf{Avoidance} (exit the activity) --- \textbf{Retention} (deductibles, uninsurable or uneconomic exposures)
  \end{itemize}
\end{itemize}
\end{frame}

\begin{frame}{Financial Risks: Market Price Volatility}
\begin{itemize}
  \item Source: price movements in organized financial markets --- FX, interest rates, commodities, equities.
  \item \alert{Financial risk $\neq$ any risk that loses money.} Confusing the two is the most common classification error in this chapter.
\end{itemize}
\medskip
\begin{exampleblock}{Example: a 90-day euro payable}
\begin{itemize}
  \item U.S.\ manufacturer owes a euro-denominated payment in 90 days.
  \item Dollar weakens from \$1.10 to \$1.20 per euro $\Rightarrow$ same obligation costs $\sim$9\% more dollars.
  \item No physical damage, no process failure, no strategic misjudgment --- a loss purely from a market price movement. That is financial risk in its most transparent form.
\end{itemize}
\end{exampleblock}
\end{frame}

\begin{frame}{What Makes Financial Risks Different}
\begin{itemize}
  \item \textbf{Two-sided}: gains as well as losses --- so management involves a deliberate hedge-or-accept choice.
  \item \textbf{Continuous}: the risk persists as long as the exposure exists, not as discrete events.
  \item \textbf{Measurable with unusual precision}: observable volatility and correlations support VaR, stress testing, sensitivity analysis.
  \item \textbf{Hedgeable}: liquid derivatives markets exist for currencies, rates, and many commodities.
  \item First management step: separate \alert{inherent} exposures (airline buying jet fuel) from \alert{incidental} ones (that airline's pension equity portfolio) --- then decide what to hedge, how much, with what.
\end{itemize}
\end{frame}

\begin{frame}{Operational Risks: The Broadest Quadrant}
\begin{itemize}
  \item Basel definition: ``the risk of loss resulting from inadequate or failed internal processes, people and systems or from external events.''
  \item Picture: the payment system crashes for six hours on a December Saturday. No asset damaged, no price moved --- but revenue disappears and customers defect.
  \item \textbf{Process} risks (designs that fail at peak volume), \textbf{people} risks (error and fraud), \textbf{technology} risks (outages, defects, \alert{cyber}), and \textbf{external events} (a supplier's factory fire, a pandemic).
  \item Generally \alert{not hedgeable} --- no derivative pays off when software fails. Tools: preventive/detective/corrective controls, business continuity planning, internal audit, and \alert{culture} --- psychological safety to raise concerns early.
\end{itemize}
\end{frame}

\begin{frame}{Business Interruption: Operational, Not Financial}
\begin{itemize}
  \item Fire shuts a plant for three months; revenue loss is real and material.
  \item But no market variable moved: prices for the product are unchanged --- \alert{the company just can't produce it}. Source: physical peril and stopped operations.
\end{itemize}
\medskip
\begin{block}{Why the category matters}
Financial risks are hedged with derivatives; there is no derivative for a production shutdown. The right toolbox is hazard controls, continuity planning, \alert{business interruption insurance}, and operational redundancy. Wrong category $\Rightarrow$ shopping in markets that don't sell what you need.
\end{block}
\end{frame}

\begin{frame}{Strategic Risks: Is the Strategy Itself Right?}
\begin{itemize}
  \item Fundamental uncertainty about markets, business model, positioning, and response to external change.
  \item Cannot be transferred or controlled away --- managed through the \alert{quality of the strategy}, monitoring and adaptation, and governance scrutiny.
  \item The hardest feature: strategic risk can originate in \emph{current strengths}. A dense store network was an advantage in 1995 and a liability a decade later. (Netflix again.)
  \item Sources: competition and substitutes, technology trajectories (ICE vs.\ electrification), regulation and politics, and \alert{reputation}.
  \item Two distinguishing traits: genuinely hard to quantify, and inseparable from strategy --- doing nothing is itself a strategic choice. The question is never \emph{whether} to accept strategic risk, but \emph{which}.
\end{itemize}
\end{frame}

\begin{frame}{Risks Cross the Quadrant Boundaries}
\begin{itemize}
  \item The taxonomy is a discipline for identification, not a filing system. The biggest enterprise risks span categories.
  \item A \alert{cybersecurity breach}:
  \begin{itemize}
    \item Operational failure (controls, unpatched vulnerability, social engineering)
    \item Hazard/liability exposure (customers, regulators seek compensation)
    \item Regulatory and strategic risk (investigations; lost differentiation if security was a selling point)
    \item Financial dimensions (investor confidence, credit standing)
  \end{itemize}
  \item Use the quadrants to ensure completeness --- then ask which risks \emph{interact}, and whether governance can see the whole picture.
  \item The goal: \alert{one integrated view, not four tidy lists}.
\end{itemize}
\end{frame}

\section{Risk Identification Methodologies}

\begin{frame}{Internal Sources}
\begin{itemize}
  \item \textbf{Historical loss data}: incidents, near-misses, control failures --- best for hazard and operational risks. Limitation: the past. Novel risks appear in no database.
  \item \textbf{Audit findings and operational metrics}: a weak control is a risk even before it's exploited; rising complaints, error rates, and turnover signal trouble early. \alert{KRIs} alongside KPIs make monitoring continuous.
  \item \textbf{Workshops, interviews, surveys}: frontline employees know the workarounds and near-misses that never reach management; facilitated cross-functional workshops force different perspectives into one conversation.
  \item \textbf{Management and expert judgment}: competitive, regulatory, and technical risks that general workshops will miss.
\end{itemize}
\end{frame}

\begin{frame}{External Sources}
\begin{itemize}
  \item \textbf{Peer Risk Factors disclosures}: what comparable firms identify as material --- surfaces risks you face but haven't articulated.
  \item \textbf{Industry loss data}: trade associations, insurers, data providers --- plausible magnitudes and peer benchmarks.
  \item \textbf{Regulatory guidance and enforcement}: enforcement against peers is risk information someone else paid for.
  \item \textbf{Environmental scanning}: technology, competition, regulation, macro, social trends --- with assigned responsibility for defined domains and periodic synthesis.
\end{itemize}
\medskip
Netflix is the model: scanned the environment, named the trend as a strategic risk, and put it in investor-facing disclosure before the disruption was obvious.
\end{frame}

\begin{frame}{Structured Techniques}
\begin{itemize}
  \item \textbf{Scenario analysis} --- narratives of how adverse situations unfold (key customer defects, plant disabled, breach, executive departure). Best for risks with no historical record and complex multi-dimensional events.
  \item \textbf{Process mapping} --- at each mapped step: what could go wrong, what controls exist, are they adequate? Productive in high-volume transaction environments.
  \item \textbf{Root cause analysis} --- ask ``why?'' until you reach drivers, not symptoms.
  \item \textbf{Bow-tie analysis} --- causes and preventive controls on one side, consequences and mitigating controls on the other; shows convergence and cascade around a single event.
\end{itemize}
\end{frame}

\begin{frame}{Making Identification Continuous}
\begin{itemize}
  \item Risk identification fails as an annual event --- risks don't wait for the workshop.
  \item Embed it where decisions are made:
  \begin{itemize}
    \item \textbf{Strategic planning}: risks of each alternative; consistency with appetite
    \item \textbf{Budgeting}: what could prevent planned performance
    \item \textbf{Projects}: identification at inception, reassessment through execution
    \item \textbf{M\&A due diligence}: strategic fit, culture, integration, target liabilities --- before capital commits
  \end{itemize}
  \item Quarterly senior-management and board risk reviews keep the \alert{risk register} reflecting current conditions, not last year's assessment.
\end{itemize}
\end{frame}

\section{The SEC Form 10-K Risk Factors Disclosure}

\begin{frame}{A Required, Structured Risk Inventory}
\begin{itemize}
  \item No other public document matches it: strategic, operational, financial, regulatory, and hazard risks in one organized section --- written by management, vetted by counsel, reviewed by the board, filed with a regulator.
  \item Legal basis: Securities Acts of 1933/1934 set the disclosure framework; risk factor requirements grew from 1970s--80s offering rules; Regulation S-K Item~503(c) (2005), reorganized as \alert{Item~105} (2019).
  \item Item 105 today: disclose \alert{material} factors that make the investment ``speculative or risky''; organize under headings; summarize if over 15 pages; no risks ``that could apply to any registrant.''
  \item The 2019 amendments were a tacit admission that many disclosures had grown so long and generic they obscured more than they revealed.
\end{itemize}
\end{frame}

\begin{frame}{What Good Disclosure Looks Like}
\begin{itemize}
  \item Typical structure: risks grouped by type, each factor headed by an adverse statement (``We depend on key suppliers, and disruption could adversely affect\ldots''), followed by why, how, and consequences.
  \item Specific beats generic:
  \begin{itemize}
    \item \emph{Weak}: ``Customer concentration risk could adversely affect our business.''
    \item \emph{Strong}: Top five customers = 45\% of revenue; one announced in-house production; a similar loss cost \$30M last year.
  \end{itemize}
  \item Content varies by industry: tech emphasizes cyber, talent, platform dependence; manufacturers emphasize supply chains and recalls; financials emphasize credit, liquidity, capital; healthcare emphasizes FDA and patents.
  \item Reading across industries with the \alert{four-quadrant lens} is one of the best ways to build identification fluency.
\end{itemize}
\end{frame}

\begin{frame}{Using Risk Factors Analytically}
\begin{itemize}
  \item \textbf{Peer inventory}: consolidate risks disclosed across the industry; check applicability; find the gaps in your own register. Collective identification effort, free of charge.
  \item \textbf{Cross-sectional comparison}: common risks define the industry (all airlines: fuel; all pharma: FDA); \alert{company-specific risks} signal particular vulnerabilities --- a concentration, a legacy platform, an investigation peers don't face.
  \item \textbf{Year-over-year changes} --- underused: newly added or expanded factors signal management has identified a new or growing concern.
  \item Caution: retrospective studies show firms often disclose risks prominently only \emph{after} they materialize. Absence from Risk Factors is not evidence of absence.
\end{itemize}
\end{frame}

\begin{frame}{Limitations: Read Critically}
\begin{itemize}
  \item \textbf{Boilerplate}: ``competition in our markets is intense and may increase'' tells you nothing company-specific.
  \item \textbf{Legal incentives}: the PSLRA (1995) safe harbor protects forward-looking statements accompanied by cautionary language --- so firms disclose broadly for protection, not clarity.
  \item \textbf{No quantification}: ``significant foreign exchange risk'' --- but what share of revenue, which currencies, what hedges, what rate move is material? Usually unstated.
  \item Supplement Risk Factors with the MD\&A, footnotes, and market risk disclosures to prioritize rather than merely catalogue.
\end{itemize}
\end{frame}

\begin{frame}{Discussion Questions}
\begin{enumerate}
  \item A three-month plant shutdown after a fire costs \$40M in lost revenue. A classmate calls this a financial risk ``because the loss is financial.'' Correct them using the source-versus-consequence distinction --- and name the right management tools.
  \item Netflix disclosed its own business model's obsolescence as a risk factor while the model was still profitable. What organizational and governance conditions make that kind of identification possible --- or impossible?
  \item Pick a cybersecurity breach at a company you know. Trace its consequences through all four quadrants. Which function should ``own'' the risk, and what does your answer imply about the taxonomy's limits?
  \item The PSLRA safe harbor rewards broad, cautious disclosure. Does that make Risk Factors more useful or less useful for risk identification? What would you change about Item~105 if you could?
\end{enumerate}
\end{frame}

\end{document}
